\section{结论}

本轮工作把备赛从“算法清单”转为可验收的模型、代码和写作链。三名成员分别承担
机理连续模型、综合优化和数据决策的首责，同时共享编程、误差分析和选题判断；
18 道真题为每条能力提供了训练入口。论文研究限定在 59 篇编号语料，59 张正式证据卡
覆盖 2892 页并可追溯，赛题和评述不混入文风样本；A/B/C 三类分别完成 20/19/20 篇
人工门禁。模板和 Skill 已形成可增量修改的独立骨架，并已吸收本轮 59 篇证据。四周计划
把尚无证据的能力转化为代码、图表、论文和复现验收。下一阶段的重点不是继续
增加算法名称，而是在限时演练中证明基线、约束、检验和表达能够稳定闭环。

\begin{thebibliography}{99}
\bibitem{papers}
2020--2025 年全国大学生数学建模竞赛 A/B/C 题编号论文语料（本地归档），
59 篇，2892 页。

\bibitem{problems}
全国大学生数学建模竞赛赛题原文：2020--2025 年 A/B/C 题（本地归档）。

\bibitem{commentaries}
全国大学生数学建模竞赛专家评述：2020--2025 年相关题目（本地归档）。

\bibitem{mindmap}
数学建模学习思维导图 \texttt{1.png}（用户提供，本地文件，核验日期：
2026-07-28）。
\end{thebibliography}

\appendix
